\documentclass[11pt]{article}
\usepackage[margin=1in]{geometry}
\usepackage{amsmath,amssymb,amsthm}
\usepackage{bm}
\usepackage{enumitem}
\usepackage[hidelinks]{hyperref}
\usepackage{booktabs}
\usepackage{graphicx}
\usepackage{caption}
\usepackage{xcolor}
\usepackage{longtable}
\usepackage{multirow}

\newcommand{\R}{\mathbb{R}}
\newcommand{\E}{\mathbb{E}}
\newcommand{\proto}{\Theta}
\newcommand{\main}{\textsc{main-full}}
\newcommand{\orig}{\textsc{original}}
\newcommand{\rid}{\mathrm{IC}_{\mathrm{rank}}}

\title{\vspace{-2.5em}\textbf{Closing the Research Loop}\\[0.3em]
\large Decay-Driven Dynamic Re-Mining:\\
Design, Evidence, and the Experiment That Tests It}
\date{28 September 2026}

\sloppy
\begin{document}
\maketitle
\vspace{-2em}

\begin{abstract}
\noindent
The system mines alpha factors once, from a human-written direction, and trades
the result indefinitely. On the first genuinely leak-free CSI300 window this does
not survive: rank IC decays from $0.1104$ to $\sim\!0.045$ within three years and
cumulative excess return over the cap-weighted, dividend-reinvested CSI300 is
$-41.3\%$ across 2017--2025. We present (i) the one intervention already measured
to work --- an annual combiner refit, which turns a $-18.8\%$ cumulative excess
into $+2.7\%$ on \emph{unchanged} factors; (ii) a design that closes the research
loop by monitoring decay and re-mining against it, with deterministic rules as a
hard floor and an LLM agent permitted only to accelerate them; and (iii) a
nine-arm historical replay that isolates every mechanism, including two non-LLM
control arms that test for temporal leakage at a layer the main ablation cannot
reach. We then argue, from metrics rather than assertion, that this work closes
the \emph{research} half of the quant loop and cannot close the \emph{trading}
half: predictive power is present and industry-normal, and the shortfall is
located entirely in portfolio construction.
\end{abstract}

\section{Where we are}

All figures below use protocol
\texttt{protocol\_csi300\_lightgbm\_topk\_decay.yaml}: fit 2005-01-01 to
2016-12-31 (purged to 2016-12-09), scored 2017-01-01 to 2025-12-31, LightGBM
combiner, \texttt{topk\_dropout} construction ($k=50$, \texttt{n\_drop}$=5$,
\texttt{max\_weight}$=0.03$), the full $\kappa$ cost model, and benchmark
\texttt{SH000300TR} --- the official cap-weighted CSI300 \emph{total-return}
index. Because \texttt{costs.net\_return} computes
$\text{gross}-\text{benchmark}-\text{charges}$, every return reported here is
\textbf{excess}, and CRR denotes \textbf{cumulative excess return}.

\begin{table}[h]
\centering
\caption{Leak-free out-of-sample baseline, 2017--2025.}
\begin{tabular}{lrrrrr}
\toprule
configuration & $\rid$ & ARR & \textbf{CRR} & IR & TC \\
\midrule
\main{} (147 factors)  & 0.0565 & $-5.96\%$ & $\bm{-41.32\%}$ & $-0.32$ & 0.37 \\
\orig{} (147 factors)  & 0.0488 & $-9.78\%$ & $-59.07\%$ & $-0.59$ & 0.33 \\
null model (no factors) & --- & $-14.33\%$ & $-73.87\%$ & $-1.05$ & --- \\
\bottomrule
\end{tabular}
\end{table}

Two properties of this baseline drive everything that follows.

\paragraph{Signal decays on a roughly three-year clock.}
Rank IC by year for \main{}: $0.1104$, $0.0898$, $0.0587$, $0.0543$, $0.0504$,
$0.0524$, $0.0160$, $0.0288$, $0.0469$. The decline is monotone in distance from
the fit window and then plateaus near $0.045$. The 2023 trough is genuine: IC goes
negative for the reduced zoo that year.

\paragraph{The returns are beta, not alpha.}
Correlation of annual excess return with the \emph{benchmark's} return is
$\bm{-0.73}$; with the book's own rank IC it is $+0.11$. Signal quality explains
about $1\%$ of the variance in excess return. Every profitable year is a year the
index fell. The decisive case is 2017: rank IC $0.1104$, the highest this system
has ever produced out of sample and above every break-even estimate, returned
$-5.45\%$ because the index rose $24.25\%$.

\section{What already works: the annual combiner refit}

The cheapest available intervention changes no factors at all. It refits the
LightGBM combiner each year on data strictly before the year being traded
(fit $\le Y-1$, report $Y$, expanding window). The factor library is frozen
throughout, so any gain is attributable to the refit alone.

\begin{table}[h]
\centering
\caption{\main{}, 147 factors: frozen 2016 fit versus annual rolling refit.
Both columns are excess over \texttt{SH000300TR}, net of $\kappa$ costs.}
\begin{tabular}{lrrrrrr}
\toprule
& \multicolumn{2}{c}{frozen (fit $\le$ 2016)} & \multicolumn{2}{c}{annual refit} & \multicolumn{2}{c}{improvement} \\
\cmidrule(lr){2-3}\cmidrule(lr){4-5}\cmidrule(lr){6-7}
year & $\rid$ & ARR & $\rid$ & ARR & $\Delta\rid$ & $\Delta$ARR \\
\midrule
2021 & 0.0504 & $+5.82\%$  & 0.0586 & $+11.40\%$ & $+0.0082$ & $+5.58$ pp \\
2022 & 0.0524 & $+1.81\%$  & 0.0598 & $+8.98\%$  & $+0.0074$ & $+7.17$ pp \\
2023 & 0.0160 & $-7.42\%$  & 0.0298 & $-1.99\%$  & $+0.0138$ & $+5.43$ pp \\
2024 & 0.0288 & $-4.31\%$  & 0.0533 & $-3.75\%$  & $+0.0245$ & $+0.56$ pp \\
2025 & 0.0469 & $-14.94\%$ & 0.0616 & $-10.25\%$ & $+0.0147$ & $+4.69$ pp \\
\midrule
mean & 0.0389 & --- & 0.0526 & --- & $+35\%$ & $\bm{+4.69}$ \textbf{pp} \\
\textbf{stitched} & & $\bm{-4.24\%}$ & & $\bm{+0.55\%}$ & & \\
\textbf{CRR} & & $\bm{-18.8\%}$ & & $\bm{+2.68\%}$ & & $\bm{+21.5}$ \textbf{pp} \\
\bottomrule
\end{tabular}
\end{table}

\noindent
Every one of the five years improves, on both rank IC and ARR. Mean rank IC rises
$35\%$; the transfer coefficient rises from $0.366$ to $0.386$. \textbf{This is the
only configuration of the nine we have tested with a positive CRR}, and it is
currently \textbf{not done at all} --- the production system fits once and never
refits. The frozen stitched figure is compounded from the per-year values; the
refit figure is computed on the concatenated daily series.

Note also that 2023 --- the year attributed in the literature to an A-share regime
rotation --- more than halves its loss under a refit, and its rank IC rises $86\%$.
Roughly half of what reads as a regime break is combiner staleness.

\paragraph{Factor freshness is worth about twice as much again.}
Fresh factors score $0.1001$ (2017--18) and $0.1136$ (2016) against $0.0526$ for
stale factors with a fresh combiner and $0.0439$ for both stale. Combining the two
through $\mathrm{IR}=\mathrm{TC}\cdot \mathrm{IC}\cdot\sqrt{N}$: both stale
$0.30$, refit-only $0.36$, re-mined and refit $\bm{0.68}$ --- the first
configuration above the conventional $0.5$ viability bar. \emph{This figure carries
a caveat that the experiment is designed to test; see \S\ref{sec:leak}.}

\section{The design}

Mining cadence is the primary lever and combiner refit the secondary one, and both
belong inside one loop. The design has a deterministic layer that always runs and
an advisory layer that may only make it act sooner.

\paragraph{Deterministic layer (the hard floor).}
Per-factor rolling 63-day IC is compared against the factor's validation IC at
admission. Tiers, on \emph{sustained} breaches rather than touches:
\emph{healthy} ($\ge 70\%$ of baseline); \emph{soft decay} ($<70\%$ for 30+ days)
$\to$ half weight; \emph{hard decay} ($<50\%$ for 60+ days) $\to$ removed from
scoring, UUID retained and re-tested quarterly. Cadences: monitor \textbf{daily};
refit \textbf{monthly} floor; re-mine \textbf{annual} floor, fired earlier when
$\ge 33\%$ of the active set is hard-decayed or the book's own rolling IC breaches
$50\%$ of baseline for 60 days.

Annual is the \emph{minimum} defensible re-mine cadence, not a cautious one: the
published estimate for a medium-frequency factor's half-life is now around 18
months, against 5--7 years pre-AI.

\paragraph{Advisory layer (the decay agent).}
Because the deterministic rule cannot fire before a 30-day sustained breach, there
is a structural 30--60 day latency to buy back. An LLM agent runs weekly and on
every tier transition, and may: trigger a refit or re-mine early, or demote a
factor early. It may delay exactly one thing --- a \emph{soft} demotion, via a
regime exemption that must cite a computed market condition, expires after one
quarter, may never cover more than $25\%$ of the active set, and may never apply
to hard decay. It may never postpone a scheduled refit or re-mine. A malformed,
timed-out, or absent agent is a non-event: the deterministic layer proceeds
unchanged.

\paragraph{Soft decay under a tree model.}
``Half weight'' is not literally implementable: gradient-boosted trees are
invariant to monotone feature scaling and the LightGBM path produces no per-factor
weights. The combiner already fits one model per seed and averages the
predictions, so soft decay is expressed as \textbf{feature-subset bagging}: the
factor's column enters only $\lceil p\cdot K\rceil$ of the $K$ seed models, with
$p\in\{1,0.5,0\}$ reproducing the tier multipliers exactly.

\paragraph{The library, and what actually trades.}
The library is append-only and unbounded: every factor ever mined, admitted or
rejected, retained under its UUID with its full trajectory. What trades is a
\emph{capped, gated} active set, selected per refit by tier, then pathology, then
an orthogonality gate ($\rho_{\max}\le\bar\rho$, so no two bets are the same ---
currently disabled by default), then capacity-with-replacement, then re-measured
marginal contribution against the set as it now stands.

\paragraph{Direction selection, without a news corpus.}
Each re-mine derives its own directions from a \textbf{computed market-state
report} over data $\le T$: cross-sectional dispersion, breadth, the size-decile
spread, volatility term structure, return autocorrelation at 1/5/20 days,
turnover, overnight-gap share, industry dispersion, and mean pairwise correlation.
No news is read. This is deliberate: the 2023 rotation \emph{is} a rise in
dispersion and an inversion of the size spread, and the agent sees those as
numbers with no label. Nothing about the regime narrative is ever stated to the
model.

\paragraph{Parent selection.}
Selection today is pure exploitation over a pool destroyed at the end of each run,
and it has produced operator monoculture twice. The library becomes a MAP-Elites
style archive binned by operator class, horizon, regime profile and decay tier,
with parents drawn on a logged budget: $50\%$ exploit (best by shrunk marginal
contribution among healthy factors), $30\%$ explore (sparse archive cells), $20\%$
resurrect (hard-decayed factors whose regime profile matches the current market
state). Selection within a pool is UCB over \emph{cells}, not factors, because
individual fitness is noisy enough that three of four verdicts flip across
combiner seeds.

\section{The experiment}

Validation is a historical replay: at each trigger point $T$, everything ---
market state, direction, mining, admission, active-set gate, combiner fit --- uses
data $\le T$ only; the following year is scored and never revisited; per-arm daily
returns are stitched into one continuous series.

\begin{table}[h]
\centering
\caption{The nine arms. ``Schedule'' distinguishes the deterministic floor alone
from the floor plus the advisory decay agent.}
\small
\begin{tabular}{llllp{5.4cm}}
\toprule
arm & directions & schedule & re-mines & what it isolates \\
\midrule
\textbf{C} & --- & none & no & static baseline; the number to beat (CRR $-41.32\%$) \\
\textbf{D} & --- & calendar & no & refit only; already measured (CRR $+2.68\%$) \\
\addlinespace
\textbf{A} & fixed & deterministic & yes & the value of re-mining itself \\
\textbf{B} & agent & deterministic & yes & the direction agent (vs A) \\
\textbf{G} & fixed & $+$ decay agent & yes & the decay agent (vs A) \\
\textbf{H} & agent & $+$ decay agent & yes & the full system; interaction of both agents \\
\textbf{H$^-$} & agent & $+$ agent, accelerate-only & yes & the regime-exemption mechanism (vs H); run only if H beats G \\
\addlinespace
\textbf{E} & Alpha158 & calendar & no & leakage control: a fixed, published, pre-LLM factor set \\
\textbf{F} & random DSL & calendar & no & leakage control: mechanically generated expressions, no survivorship \\
\bottomrule
\end{tabular}
\end{table}

\noindent
Arms A, B, G, H form a $2\times2$ factorial over \{fixed, agent\} directions
$\times$ \{deterministic, $+$agent\} scheduling. The contrasts of interest are
$A-D$ (re-mining), $B-A$ (direction agent), $G-A$ (decay agent), $H-G$ and $H-B$
(each agent in the presence of the other), and the interaction
$(H-G)-(B-A)$.

\paragraph{Metrics.}
Primary: stitched \textbf{CRR}. Secondary, per year and per arm: ARR, IR, MDD,
rank IC, IC, ICIR, TC, turnover, cost in bps, active-set size, admissions, and the
decay tier census. The benchmark's own return is reported alongside every arm,
because excess return here is beta-driven and an arm that ``wins'' in a falling
market has not necessarily produced alpha. Deflated Sharpe is computed against the
\emph{true cross-mine trial count}, which grows with every mine and correctly
raises the bar.

\paragraph{Cost.}
Annual triggers across 2017--2025 give nine points per re-mining arm. Arms G and H
share mines with A and B wherever their trigger dates coincide, so the incremental
cost of the decay-agent arms is only the divergences; the estimate is 22--26 mines
in total. Arms C, D, E and F require no mining at all. The per-mine wall-clock and
LLM cost is measured before anything is launched, and a reduced three-point trigger
set is the pre-committed fallback.

\subsection{Temporal leakage, and why two control arms}
\label{sec:leak}

The generating model's weights contain the future; no prompt retracts that. Three
layers bound what leakage can do --- context hygiene (only the train window is
disclosed, clamped to the day before the earliest fold's validation day, failing
closed), an output validator rejecting dates, tickers and named events, and a
specificity cap so a direction names mechanisms rather than answers. But bounds
are not measurements, so the experiment measures two distinct exposures.

\paragraph{Late-sample leakage: already refuted.}
If hindsight had been exploited, the factors would hold up in the leaked period.
They do not: rank IC falls monotonically with distance from the fit and troughs at
$0.0160$ in \textbf{2023}, the most-documented regime break in the sample and
precisely where a leaking model's knowledge is densest. There is also a structural
reason --- admission selects on data $\le T$, so a factor whose edge exists only in
the future is rejected before entering the library.

\paragraph{First-year leakage: \emph{not} refuted, and it matters.}
Within the pool that works before $T$, hindsight can still select those the model
knows persist longest. That inflates \textbf{year one} and converges to the honest
baseline as even the leaked picks decay --- reproducing the observed curve exactly.
The 2023 trough does not exclude it, since by 2023 a leaked pick is dead too. And
the $A$ versus $B$ ablation is \textbf{blind} to it, because both arms author
factors with the same model: this leakage is common mode and differences out.

This bears directly on the project's central justification. The $+90\%$
factor-freshness figure, which carries re-mined-and-refit to gross IR $0.68$, is a
\emph{first-year} measurement --- and an annually re-mining loop lives on exactly
that number.

Hence arms \textbf{E} and \textbf{F}, which contain no LLM in generation and
therefore cannot leak. Arm E is Alpha158, a fixed published pre-LLM set; arm F is
mechanically generated DSL expressions, which arm E cannot replace because
Alpha158 is a \emph{survivor} set and may show unusually flat decay for reasons
unrelated to leakage. The test compares the \textbf{ratio} $\rid^{\text{year-1}} /
\rid^{\text{steady-state}}$ (a ratio, because the controls are weaker in level). A
materially larger ratio for the mined arms is the leakage signature, isolated at
the layer the main ablation cannot reach. \textbf{Until these arms clear year one,
the deployable expectation is the steady-state rank IC ($\sim\!0.045$--$0.05$), not
the first-year value ($\sim\!0.11$), and $+90\%$ is quoted as an upper bound.}

\section{What this closes, and what it cannot}

\subsection{The research half is present and industry-normal}

\begin{table}[h]
\centering
\caption{Predictive-power metrics. These are construction-, cost- and
benchmark-independent: they are properties of the signal alone.}
\begin{tabular}{lrrl}
\toprule
metric & full 2017--2025 & 2017 & typical usable range \\
\midrule
IC            & 0.0425 & 0.0941 & 0.02--0.05 \\
Rank IC       & 0.0565 & 0.1104 & 0.03--0.05 \\
ICIR          & 0.315  & 0.743  & 0.3--0.5 \\
Rank ICIR     & 0.409  & 0.845  & 0.3--0.5 \\
\bottomrule
\end{tabular}
\end{table}

\noindent
On leak-free out-of-sample data the mined library sits inside the normal range for
a usable equity signal on every predictive-power metric, and its first out-of-sample
year is comfortably above it. \textbf{The system finds alpha.} That is the research
question, and it is answered affirmatively.

\subsection{The trading half is where it is lost, and the metrics locate it exactly}

Write the standard decomposition
$\mathrm{IR}=\mathrm{TC}\cdot\mathrm{IC}\cdot\sqrt{N}$. With
$\mathrm{IC}=0.0425$, $N\approx 300$ so $\sqrt{N}=17.3$, and
$\mathrm{TC}=0.37$:
\[
\mathrm{IR}_{\text{gross}} \;=\; 0.37 \times 0.0425 \times 17.3 \;\approx\; 0.27\text{--}0.36 .
\]
Every term is individually normal. \textbf{A transfer coefficient of $0.37$ is
squarely inside the industry range of $0.3$--$0.6$ for a long-only book} --- there
is no anomalous loss to recover. The product is sub-viable because of the
\emph{mandate}: $N\approx 300$ names, $k=50$ holdings, long-only, benchmarked to a
cap-weighted index with no hedging instrument. Five independent measurements
confirm the constraint is construction, not signal:

\begin{enumerate}[leftmargin=*,itemsep=2pt]
\item \textbf{The null model reproduces the pathology.} A book holding \emph{no}
mined factors shows correlation $-0.64$ between excess return and the benchmark,
against $-0.73$ for \main{}. The beta shortfall is a property of the construction,
not of the factors.
\item \textbf{The position cap is not binding.} Raising \texttt{max\_weight} from
$0.03$ to $0.10$ left TC at $0.36$ and cost $10.5$ pp. The mean-variance optimiser
\emph{diversified} from 100 to 134 names rather than concentrating --- a
total-risk optimiser answering a benchmark-relative mandate.
\item \textbf{Relaxing long-only raises TC and still loses.} Signed (dollar-neutral)
top-$k$ lifted TC from $0.37$ to $\bm{0.51}$ at an unchanged cap --- confirming
long-only, not the cap, is the binding constraint --- yet absolute long/short
return was $-1.65\%$/yr with 2 of 9 positive years.
\item \textbf{Three of nine tested configurations are worse than holding no factors
at all}, including Kelly/mean-variance at CRR $-78.95\%$ against the null model's
$-73.87\%$.
\item \textbf{The best signal year is a losing year.} 2017, rank IC $0.1104$, ICIR
$0.743$: $-5.45\%$, because the index rose $24.25\%$.
\end{enumerate}

\subsection{The asymmetry, stated plainly}

\begin{table}[h]
\centering
\caption{What each half of the loop actually has.}
\small
\begin{tabular}{p{6.4cm}p{6.4cm}}
\toprule
\textbf{Factor side (research)} & \textbf{Portfolio side (trading)} \\
\midrule
LLM search loop with mutation, crossover and refinement &
\texttt{topk}$=50$ \\
Admission gate on marginal contribution tested against its own noise &
\texttt{max\_weight}$=0.03$ \\
Benjamini--Hochberg FDR control; deflated Sharpe against trial count &
\texttt{n\_drop}$=5$ \\
Orthogonality gate, diversity scoring, lineage and verdict tracking & \\
Per-factor tear sheets, diagnosis and routing & \\
\midrule
\emph{A research system.} & \emph{Three constants, fixed once by a 54-config sweep on a single 2021 window.} \\
\bottomrule
\end{tabular}
\end{table}

\noindent
The objective the search maximises is a function of IC. \textbf{Nothing in the
system optimises IR or CRR.} The re-mining design in this document makes the
research half close properly --- it monitors its own decay, re-mines against it,
and refits continuously. It will raise IC, and on the measured decomposition that
raises gross IR from $0.30$ to perhaps $0.68$.

It will not fix a transfer coefficient that is already normal, a benchmark that
cannot be tracked long-only, or a construction stage with no search loop, no
admission criterion and no multiple-testing control. \textbf{Converting demonstrated
predictive power into realised profit requires a second system, built for
construction the way this one was built for factors} --- searching over
constraints, hedges, weighting schemes and rebalancing rules, with an objective
that is the realised net IR rather than the signal's IC.

\paragraph{The falsifiable form of this claim.} If the conclusion is wrong, the
replay will say so: an arm that materially improves IC should then also improve
CRR. The measured relationship says it will not --- correlation between annual
excess return and rank IC is $+0.11$, against $-0.73$ with the benchmark. We will
report whichever answer the data gives.

\section{What is not tested here}

Shorting and borrow costs (deferred to the construction phase, and the signed
result above says it is not a quick win); live news; real capital and a broker
integration; and any change to the construction stage itself. NAV is simulated and
compounds from \textyen100M, which at that scale leaves the market-impact term
effectively inert --- it matters as the capital base changes, not now.

\end{document}
